\pdfinfoomitdate=1
\pdftrailerid{}
\pdfsuppressptexinfo=15
\documentclass[11pt]{article}
\usepackage[T1]{fontenc}
\usepackage{lmodern}
\usepackage{amsmath,amssymb,amsthm}
\usepackage[hidelinks]{hyperref}
\usepackage{microtype}
\usepackage[margin=1in]{geometry}

\newtheorem{lemma}{Lemma}
\newtheorem{remark}[lemma]{Remark}
\newcommand{\MC}{\operatorname{MC}}
\newcommand{\ad}{\operatorname{ad}}
\newcommand{\im}{\operatorname{im}}
\newcommand{\QQ}{\mathbf Q}

\title{A fixed-two-jet extension lemma for complete strict dg Lie algebras}
\author{}
\date{}

\begin{document}
\maketitle

We use the ordinary cohomological convention.  Thus a dg Lie algebra
$(\mathfrak h,\partial,[-,-])$ has $\partial$ of degree $1$, its bracket has
degree $0$, and Maurer--Cartan elements have degree $1$.  No transferred or
higher operations occur in this note.

Give
\[
 s\mathfrak h[[s]]=\prod_{j\geq 1}s^j\mathfrak h
\]
the complete separated filtration
$F^q=s^q\mathfrak h[[s]]$.  Its Maurer--Cartan equation is
\[
 \partial\alpha+\frac12[\alpha,\alpha]=0.
\]

\begin{lemma}[fixed-two-jet strict extension]\label{lem:strict-extension}
Let $\mathfrak h$ be a dg Lie algebra over a field of characteristic zero.
Suppose
\[
 \bar\alpha=s a_1+s^2a_2+s^3a_3
       \in\MC\bigl(\mathfrak h\otimes (s)/(s^4)\bigr),
 \qquad a=a_1.
\]
If
\[
 \ker\bigl(\ad_a:H^2(\mathfrak h)\longrightarrow H^3(\mathfrak h)\bigr)
 =
 \im\bigl(\ad_a:H^1(\mathfrak h)\longrightarrow H^2(\mathfrak h)\bigr),
 \tag{1}
\]
then there is
\[
 \alpha_\infty\in\MC(s\mathfrak h[[s]])
 \quad\text{such that}\quad
 \alpha_\infty\equiv\bar\alpha\pmod {s^3}.
\]
In particular, the coefficients of $s$ and $s^2$ are unchanged.
\end{lemma}

\begin{proof}
For $\gamma\in s\mathfrak h^1[[s]]$, put
\[
 \mathcal F(\gamma)=\partial\gamma+\frac12[\gamma,\gamma].
\]
The dg Lie identities give the strict Bianchi identity
\[
 \partial\mathcal F(\gamma)+[\gamma,\mathcal F(\gamma)]=0.
 \tag{2}
\]
Indeed, since $|\gamma|=1$,
\[
 \partial\left(\frac12[\gamma,\gamma]\right)
 =\frac12\bigl([\partial\gamma,\gamma]-[\gamma,\partial\gamma]\bigr)
 =-[\gamma,\partial\gamma],
\]
and graded Jacobi in characteristic zero gives
$[\gamma,[\gamma,\gamma]]=0$.

Choose any lift $\gamma$ of $\bar\alpha$ to
$s\mathfrak h^1[[s]]$.  Its curvature is $O(s^4)$.  Suppose inductively that
for some $N\geq4$ we have
\[
 \mathcal F(\gamma)
 =s^N r_N+s^{N+1}r_{N+1}+O(s^{N+2}),
 \qquad \gamma=sa+O(s^2).
 \tag{3}
\]
The coefficients of $s^N$ and $s^{N+1}$ in (2) are respectively
\[
 \partial r_N=0,
 \qquad
 \partial r_{N+1}+[a,r_N]=0.
 \tag{4}
\]
Thus $r_N$ is a degree-two cocycle and
$[a,r_N]$ is a boundary.  Equivalently,
\[
 [r_N]\in
 \ker\bigl(\ad_a:H^2(\mathfrak h)\longrightarrow H^3(\mathfrak h)\bigr).
\]
By (1), there is a closed $v\in\mathfrak h^1$ for which
\[
 [r_N]+[a,v]=0\quad\text{in }H^2(\mathfrak h).
\]
Choose $w\in\mathfrak h^1$ satisfying
\[
 r_N+[a,v]=\partial w.                                      \tag{5}
\]
Set
\[
 \delta=s^{N-1}v-s^Nw,
 \qquad \gamma'=\gamma+\delta.                              \tag{6}
\]
The exact curvature-change formula is
\[
 \mathcal F(\gamma+\delta)-\mathcal F(\gamma)
 =\partial\delta+[\gamma,\delta]+\frac12[\delta,\delta].   \tag{7}
\]
Here $\partial v=0$, $\gamma-sa=O(s^2)$, and $2N-2\geq N+1$.
Consequently the coefficient of $s^N$ on the right of (7) is
\[
 [a,v]-\partial w.
\]
Together with (5), this cancels $r_N$, while every omitted term has order at
least $N+1$.  Hence
\[
 \mathcal F(\gamma')=O(s^{N+1}).
\]

Iterating (6) gives a Cauchy sequence.  More explicitly, the step indexed by
$N$ changes only coefficients in orders at least $N-1$, so every fixed
coefficient eventually stabilizes.  Completeness gives a coefficientwise
limit $\alpha_\infty$, continuity of $\partial$ and the bracket gives
\(
 \mathcal F(\alpha_\infty)\in\bigcap_Ns^N\mathfrak h[[s]]=0,
\)
and separatedness is used in the last equality.  Since the first correction
has order $s^3$, the limit agrees with $\bar\alpha$ modulo $s^3$.
\end{proof}

\begin{remark}[the initial $s^4$ step]
If the chosen lift has no terms beyond $s^3$, its first curvature
coefficients are
\[
 r_4=[a_1,a_3]+\frac12[a_2,a_2],
 \qquad r_5=[a_2,a_3].
\]
Equation (4) says
\[
 \partial r_4=0,
 \qquad \partial r_5+[a_1,r_4]=0.
\]
The correction is $s^3v_3-s^4w_4$.  This displays directly why the first and
second-order coefficients are fixed.
\end{remark}

\begin{remark}[what is not preserved]
The conclusion is deliberately congruence modulo $s^3$, not modulo $s^4$.
At the first step the term $s^3v_3$ may change $a_3$.  To retain the entire
given element modulo $s^4$, one would need the stronger condition
$[r_4]=0$ in $H^2(\mathfrak h)$, so that a boundary correction beginning in
order $s^4$ suffices.  Exactness (1) alone does not imply this stronger
vanishing.
\end{remark}

\end{document}
